% ============================================================================
%  GENERALIZING the knowability--capacity threshold beyond 1-D.
%  Companion to sections/knowability_capacity.tex (the proven 1-D result).
%
%  Each result is explicitly labelled:
%     (PROVEN)            -- complete elementary proof, no extra hypothesis.
%     (PROVEN, COND.)     -- proven UNDER a stated structural condition.
%     (NUMERICAL)         -- validated by Monte-Carlo only (stated as such).
%     (CONJECTURE/OPEN)   -- precise conjecture with the exact missing piece.
%
%  All numbers validated by
%     experiments/kbound/theory_validation/val_knowability_capacity_general.py
%     -> results_knowability_capacity_general.json   (fixed seed 20260619).
%
%  HONEST HEADLINE (see Remark rmk:gen-scope): this reaches a STRONG-BUT-BOUNDED
%  generalization, NOT a paradigm-level general capacity. Stage 1 (multivariate
%  Gaussian-location) lifts EXACTLY for the aligned concept class and is fully
%  characterized (incl. its breaker) for the free halfspace class. Stage 2
%  (exponential / MLR family) gives a matched threshold UNDER an MLR+log-concavity
%  condition, with an explicit non-log-concave breaker. Stage 3 proves the most
%  general (distribution-free) CONVERSE and isolates exactly why a single
%  computable achievability threshold needs a regularity/identifiability condition.
% ============================================================================
\section{Toward a general knowability--capacity: how far the threshold lifts}
\label{sec:knowcap-gen}

The $1$-D theorem (\S\ref{sec:knowcap}) gives a genuine capacity $\tau=1$ for
sign-of-benefit in the Gaussian-location two-point model. We now ask how much of
that survives in higher dimension, in non-Gaussian location families, and
distribution-free. The answer is a clean layered one: the converse is essentially
universal, the achievability is exact under an alignment/MLR regularity condition,
and a single \emph{computable} capacity provably needs such a condition. We label
every claim and exhibit the breakers that delimit each regime.

% ----------------------------------------------------------------------------
\subsection{Stage 1: multivariate Gaussian-location}\label{sec:gen-mvg}
% ----------------------------------------------------------------------------
\paragraph{Model.} $P_X=\mathcal N(0,\Sigma)$, $Q_X=\mathcal N(\mu,\Sigma)$ with
$\Sigma\succ0$ known and $\mu$ unknown (estimated from an unlabeled target sample).
The deployed rule is the halfspace $g=f_0=\mathbf 1[w^\top x>0]$ and the shipped
candidate is the \emph{parallel-shifted} halfspace $f_a=\mathbf 1[w^\top x>a]$,
$a>0$ known. The disagreement region is the \emph{slab}
$D=\{x:0<w^\top x<a\}$. Concepts are halfspaces
$\eta_{v,\theta}(x)=\mathbf 1[v^\top x>\theta]$ (the label mechanism), with
$\theta_S$ the source-calibrated offset along a fixed direction. Write the
\emph{whitened (Mahalanobis) projection onto the rule normal}
\begin{equation}\label{eq:gen-proj}
  s \;:=\; w^\top x,\qquad s\sim_Q \mathcal N(m,\sigma^2),\quad
  m=w^\top\mu,\ \ \sigma^2=w^\top\Sigma w,
\end{equation}
and for a concept direction $v$ the second coordinate $u:=v^\top x$, so that
$(s,u)$ is bivariate Gaussian under $Q$ with correlation (the \emph{Mahalanobis
cosine} between $w$ and $v$)
\begin{equation}\label{eq:gen-rho}
  \rho \;=\; \rho(w,v)\;=\;\frac{w^\top\Sigma v}{\sqrt{w^\top\Sigma w}\,\sqrt{v^\top\Sigma v}}
  \;=\;\cos\angle_\Sigma(w,v).
\end{equation}
Because $f_0,f_a$ and $\eta_{v,\theta}$ depend on $x$ only through $(s,u)$, the
entire benefit is a functional of this bivariate Gaussian:
\begin{equation}\label{eq:gen-delta}
  \Delta \;=\; R_Q(f_0)-R_Q(f_a)\;=\;\mathbb E_Q\big[\mathbf 1_D\,(1-2\eta_{v,\theta})\big]
        \;=\; Q(D)\;-\;2\,Q\big(D\cap\{u>\theta\}\big).
\end{equation}

\begin{lemma}[Whitened reduction; PROVEN]\label{lem:gen-reduce}
$\Delta$ in \eqref{eq:gen-delta} depends on the dimension $d$, on $\mu$ and on
$\Sigma$ only through the bivariate law of $(s,u)$ under $Q$, i.e.\ through the four
scalars $(m,\sigma^2,\ v^\top\mu,\ v^\top\Sigma v)$ and the correlation $\rho$
of \eqref{eq:gen-rho}. In particular the deployed and candidate rules depend on $x$
only through $s$, whose $Q$-law is the one-dimensional Gaussian
$\mathcal N(w^\top\mu,\,w^\top\Sigma w)$.
\end{lemma}
\begin{proof}
$f_0=\mathbf 1[s>0]$, $f_a=\mathbf 1[s>a]$, $\eta_{v,\theta}=\mathbf 1[u>\theta]$ are
$\sigma(s,u)$-measurable, and the $0/1$ excess loss of $f_0$ over $f_a$ is supported
on $D=\{0<s<a\}$ where it equals $1-2\eta_{v,\theta}$ (exactly as in
Lemma~\ref{lem:knowcap-identity}: on $D$, $f_0{=}1,f_a{=}0$, and the conditional mean
of $\mathbf 1[Y{=}0]-\mathbf 1[Y{=}1]$ is $1-2\eta$). Taking $Q$-expectation gives
\eqref{eq:gen-delta}, a functional of the $Q$-law of $(s,u)$. A non-degenerate affine
image of a Gaussian is Gaussian with the stated moments.
\end{proof}

\paragraph{(a) Aligned concept class --- exact lift.}
Take the concept normal \emph{parallel to the rule normal}, $v=w$ (the concept
``reads the same direction'' the deployed rule acts on). Then $u=s$, the problem
collapses onto the single whitened line $\mathrm{span}(w)$, and \eqref{eq:gen-delta}
becomes literally the $1$-D identity of Lemma~\ref{lem:knowcap-identity} in the
standardized variable $z=(s-m)/\sigma$. Define the admissible aligned class and the
lifted invariant in Mahalanobis units:
\begin{equation}\label{eq:gen-Cpar}
  \mathcal C_\varepsilon^{\parallel}=\Big\{\theta_T:\ \big|\Phi(\tfrac{\theta_T-m}{\sigma})
  -\Phi(\tfrac{\theta_S-m}{\sigma})\big|\le\varepsilon\Big\},\quad
  K_\parallel=\frac{\big|\tfrac12[\Phi(\tfrac{-m}{\sigma})+\Phi(\tfrac{a-m}{\sigma})]
  -\Phi(\tfrac{\theta_S-m}{\sigma})\big|}{\varepsilon}.
\end{equation}

\begin{theorem}[Stage 1, aligned: exact Mahalanobis lift; PROVEN]\label{thm:gen-aligned}
With $v=w$ and $Q_X$ known, the population knowability--capacity
Theorem~\ref{thm:knowcap} holds verbatim with $(\mu,a,\theta_S)$ replaced by their
whitened images $(m,\sigma)$ in \eqref{eq:gen-proj}: with $\tau=1$,
$K_\parallel>1$ $\Rightarrow$ every admissible aligned concept has
$\sign\Delta=\sign(\theta_S-\theta_0)$ (achievability, error $0$); $K_\parallel<1$
$\Rightarrow$ two admissible aligned concepts share $Q_X$ and have opposite
$\sign\Delta$, so the minimax label-free error is $\tfrac12$. The flip locus is the
$Q$-median of the slab along $w$, $\theta_0=m+\sigma\,\Phi^{-1}\!\big(\tfrac12[\Phi(\tfrac{-m}{\sigma})+\Phi(\tfrac{a-m}{\sigma})]\big)$.
The finite-$n$ Le Cam boundary layer (Proposition~\ref{prop:knowcap-finiten}) carries
over with $\mu$ replaced by $m=w^\top\mu$ and the noise scale by $\sigma/\sqrt n$,
since $\bar s=w^\top\bar x\sim\mathcal N(m,\sigma^2/n)$ is sufficient for $m$.
\end{theorem}
\begin{proof}
By Lemma~\ref{lem:gen-reduce} with $v=w$, $\Delta$ is the $1$-D benefit of
Lemma~\ref{lem:knowcap-identity} for the standardized location problem
$z\sim\mathcal N(m/\sigma,1)$ with boundaries $0,a,\theta$ rescaled by $\sigma$;
\eqref{eq:gen-Cpar} is exactly \eqref{eq:knowcap-class} in $z$. Apply
Theorem~\ref{thm:knowcap}. The converse two worlds again share the \emph{same}
$Q_X=\mathcal N(\mu,\Sigma)$ (the concept lives in the unobserved $P(Y\mid X)$
channel), so $\TV=0$ for every unlabeled statistic and every $n$. The sufficiency of
$\bar s$ for $m$ is standard. \emph{(Validated: $20{,}000$ random draws in
$d\in\{2,3,4\}$, equivalence $K_\parallel>1\iff$ identifiable with $0$ mismatches off
a $10^{-6}$ shell; \texttt{val\_knowability\_capacity\_general.py} Stage~1a.)}
\end{proof}

\paragraph{(b) Free halfspace concept class --- exact obstruction and its capacity.}
If the concept normal $v$ is \emph{not} constrained to be parallel to $w$, the flip
locus moves. For a fixed $v$, $\Delta(\theta)$ in \eqref{eq:gen-delta} is still
nondecreasing in $\theta$ with a single root $\theta_0(v)$ characterized by
$Q\big(D\cap\{u>\theta_0\}\big)=\tfrac12 Q(D)$: \emph{$\theta_0(v)$ is the
$Q$-conditional median of $u=v^\top x$ given $x\in D$.} Its level in the concept's
own Mahalanobis mass coordinate,
\begin{equation}\label{eq:gen-pmed}
  p_{\mathrm{med}}(v)\;:=\;\Phi\!\Big(\tfrac{\theta_0(v)-v^\top\mu}{\sqrt{v^\top\Sigma v}}\Big),
\end{equation}
interpolates between the two extremes set by $\rho=\rho(w,v)$. In standardized
coordinates write $z=(s-m)/\sigma\sim\mathcal N(0,1)$, the slab as $z\in(\alpha,\beta)$
with $\alpha=-m/\sigma,\ \beta=(a-m)/\sigma$, and
$t=(u-v^\top\mu)/\sqrt{v^\top\Sigma v}=\rho\,z+\sqrt{1-\rho^2}\,\omega$ with
$\omega\perp z$ standard normal. Then $p_{\mathrm{med}}(v)=\Phi(\operatorname{med}(t\mid z\in(\alpha,\beta)))$ and
\begin{equation}\label{eq:gen-pmed-range}
  \rho=\pm1\ (v\parallel w):\ p_{\mathrm{med}}=
  \tfrac12[\Phi(\alpha)+\Phi(\beta)]\ (\text{band median, since }t=\pm z),
  \qquad
  \rho=0\ (v\perp_\Sigma w):\ p_{\mathrm{med}}=\tfrac12.
\end{equation}
Both endpoints are \textbf{proven exactly}. The $\rho=0$ case is structural: when $v$ is
$\Sigma$-orthogonal to $w$, $t=\omega$ is independent of $z$, so conditioning on the
slab does not move its law and $\operatorname{med}(t\mid z\in(\alpha,\beta))=0$, i.e.\
$p_{\mathrm{med}}=\tfrac12$. The $\rho=\pm1$ case is immediate since then $t=\pm z$. That
$p_{\mathrm{med}}$ is \emph{monotone} in $|\rho|$ between these endpoints (so the
achievable flip range is exactly the interval between $\tfrac12$ and the band median) is
\textbf{verified numerically} (Lemma~\ref{lem:gen-pmed-mono}, stated as a numerically
confirmed monotone-interpolation fact with a proof sketch); the \emph{worst case}
relevant to the converse --- mass $\tfrac12$, reached at $\rho=0$ --- is proven outright.

\begin{lemma}[Monotone flip interpolation; NUMERICAL, proof sketch]\label{lem:gen-pmed-mono}
With the notation above, $p_{\mathrm{med}}(v)$ is monotone in $|\rho|$, sweeping the full
interval between $\tfrac12$ (at $\rho=0$) and the band median (at $|\rho|=1$).
\emph{Sketch:} the conditional median solves
$\mathbb E[\sign(t-m_t)\mid z\in(\alpha,\beta)]=0$; since the slab's $z$-median $z^\star$
has a fixed sign relative to $0$, increasing $\rho$ tilts $t=\rho z+\sqrt{1-\rho^2}\,\omega$
monotonically toward $\pm z$, dragging $m_t$ monotonically from $0$ toward $z^\star$ by a
sign-coherent shift of the conditional law. A full proof would verify single-crossing of
the conditional-median equation in $\rho$. \emph{(Validated: $p_{\mathrm{med}}$ rises
monotonically $0.50\to$ band median as $\rho:0\to1$ across four slabs;
\texttt{val\_knowability\_capacity\_general.py} Stage~1b.)}
\end{lemma}

\begin{theorem}[Stage 1, free halfspace: the obstruction; converse PROVEN, exact $K_{\mathrm{free}}$ CONDITIONAL]\label{thm:gen-free}
Let $p_S=\Phi(\tfrac{\theta_S-m}{\sigma})$ be the calibrated boundary mass along $w$
and let the admissible free class be the halfspaces with own-normal boundary-mass
drift $\le\varepsilon$ and arbitrary direction $v$. Then the aligned invariant
$K_\parallel$ of \eqref{eq:gen-Cpar} is \emph{not} a capacity for this class. Precisely:
\begin{enumerate}\itemsep1pt
\item[\textnormal{(i)}] \textbf{(Breaker / converse; PROVEN.)} Whenever $p_S$ lies
strictly between $\tfrac12$ and $b:=\tfrac12[\Phi(\tfrac{-m}{\sigma})+\Phi(\tfrac{a-m}{\sigma})]$,
there is a regime $K_\parallel>1$ in which the aligned theorem predicts identifiability
but a tilted admissible concept (own-normal boundary mass $=p_S$, taken with
$\rho\to0$, which reaches flip mass $\tfrac12$) produces the \emph{opposite} benefit
sign. The two concepts share $Q_X$, so $\TV=0$ and the minimax error is $\tfrac12$: the
aligned $K_\parallel$ is therefore \emph{not} a capacity for the free class.
\item[\textnormal{(ii)}] \textbf{(Exact free capacity; CONDITIONAL on
Lemma~\ref{lem:gen-pmed-mono}.)} Under the monotone interpolation, the achievable flip
range is exactly the interval between $\tfrac12$ and $b$, so $\sign\Delta$ is identifiable
over the free class iff $[p_S-\varepsilon,p_S+\varepsilon]$ misses it; the governing
capacity is
\[
  K_{\mathrm{free}}\;=\;\frac{\operatorname{dist}\!\big(p_S,\ [\,\tfrac12\wedge b,\ \tfrac12\vee b\,]\big)}{\varepsilon},
  \qquad \tau=1 .
\]
\end{enumerate}
\end{theorem}
\begin{proof}
For fixed $v$, \eqref{eq:gen-delta} is nondecreasing in $\theta$ (raising $\theta$
only shrinks $\{u>\theta\}$), with unique root the conditional median $\theta_0(v)$;
hence $\sign\Delta=\sign(\theta-\theta_0(v))$. \emph{Part (i)} uses only the proven
endpoint $\rho=0\Rightarrow p_{\mathrm{med}}=\tfrac12$ of \eqref{eq:gen-pmed-range}: for
$p_S$ strictly on the $\tfrac12$-side of the aligned flip $b$, take $v$ with $\rho\to0$
and own-normal boundary mass $=p_S$; its flip sits at $\tfrac12$, on the opposite side of
$p_S$ from $b$, so $\sign\Delta$ is reversed relative to the aligned prediction. The two
concepts share $Q_X$, giving $\TV=0$ and minimax error $\tfrac12$ (Le Cam); thus
$K_\parallel$ is not a capacity. \emph{Part (ii)} additionally invokes
Lemma~\ref{lem:gen-pmed-mono} (numerically confirmed): the achievable flip levels then
fill exactly the interval between $\tfrac12$ and $b$, so a sign-flip witness exists iff
the admissible mass band $[p_S-\varepsilon,p_S+\varepsilon]$ meets that interval, i.e.\
iff $K_{\mathrm{free}}<1$. \emph{(Validated:
$p_{\mathrm{med}}$ vs $|\rho|$ traced from $0.5005$ at $\rho{=}0$ to the band median
$0.4226$ at $\rho{=}1$; an explicit breaker with $K_\parallel=1.40$ where the aligned
prediction is $\Delta=+0.14$ but a worst-case tilt of the same own-normal boundary
mass yields $\Delta=-0.14$ (opposite sign);
\texttt{val\_knowability\_capacity\_general.py} Stage~1b.)}
\end{proof}

\begin{remark}[What Stage 1 establishes]\label{rmk:gen-stage1}
The $1$-D capacity lifts to $d$ dimensions \emph{exactly} in Mahalanobis geometry
\emph{when the concept is aligned with the rule normal} (Thm~\ref{thm:gen-aligned}):
$K$ and $\tau=1$ are intact, with $\theta_0$ the slab's $Q$-median along $w$. The lift
is therefore not vacuous --- it is a genuine $d$-dimensional capacity. But the
unconstrained halfspace concept class is \emph{strictly harder}: an adversary can tilt
the concept normal toward $\Sigma$-orthogonality with $w$, dragging the flip locus to
the unconditional median (mass $\tfrac12$) and defeating any margin measured along
$w$. The honest Stage-1 statement is a capacity $K_{\mathrm{free}}$ that prices in the
worst-case tilt (Thm~\ref{thm:gen-free}); the clean closed-form $\tau=1$ survives, but
the invariant is the distance to mass $\tfrac12$, not to the band median, once the
concept direction is free.
\end{remark}

% ----------------------------------------------------------------------------
\subsection{Stage 2: MLR / log-concave location families}\label{sec:gen-mlr}
% ----------------------------------------------------------------------------
We now leave the Gaussian and ask which \emph{non-Gaussian} location families keep
the matched threshold.

\emph{Scope note.} Here ``exponential family'' abbreviates one-parameter \emph{location}
families with monotone likelihood ratio (Gaussian, Laplace, logistic location). Genuine
\emph{scale}/natural-parameter exponential families---e.g.\ the exponential or Gamma rate
family---are not location families; their capacity is not governed by the mass-coordinate
threshold below and falls under Conjecture~\ref{conj:gen-capacity}.

Let $Q_X$ have a density $q_0(x-\mu)$ for a fixed base density
$q_0$ with continuous CDF $F$; deployed $f_0=\mathbf 1[x>0]$, candidate
$f_a=\mathbf 1[x>a]$, hard-threshold concept $\eta_\theta=\mathbf 1[x>\theta]$. The
benefit identity is \emph{family-agnostic}: exactly as in
Lemma~\ref{lem:knowcap-identity} (the proof used only that on $D$ the conditional
contrast is $1-2\eta$ and that $Q$ has a CDF),
\begin{equation}\label{eq:gen-fam-delta}
  \Delta(\theta)=2F(\theta'-\mu)-F(-\mu)-F(a-\mu),\qquad\theta'=\operatorname{clip}(\theta,0,a),
\end{equation}
and the admissible class / invariant are \eqref{eq:knowcap-class},
Def.~\ref{def:knowcap-K} with $\Phi$ replaced by $F$. So far \emph{no} condition is
needed and $\Delta$ is monotone in $\theta$ for \emph{any} CDF $F$.

\begin{theorem}[Stage 2, MLR/log-concave: matched $\tau=1$; PROVEN UNDER CONDITION]
\label{thm:gen-mlr}
Assume the base family is a one-parameter \emph{location} family with strictly
\textbf{monotone likelihood ratio} in the location (equivalently here, $q_0$ strictly
positive and log-concave, hence unimodal). Then with the mass-coordinate drift class
\eqref{eq:knowcap-class}$_F$ and $\tau=1$, the population knowability--capacity holds
verbatim: $K>1\Rightarrow$ every admissible concept shares $\sign\Delta=\sign(\theta_S-\theta_0)$
(achievability), $K<1\Rightarrow$ two admissible concepts share $Q_X$ with opposite
$\sign\Delta$ (converse, $\TV=0$). Moreover the flip locus
$\theta_0(\mu)=\mu+F^{-1}\!\big(\tfrac12[F(-\mu)+F(a-\mu)]\big)$ is the band's $Q$-median
and is \textbf{strictly increasing in the unknown shift} $\mu$, so a single label-free
estimate $\hat\mu$ pins the flip side through one threshold and $K$ is a genuine scalar
capacity (the finite-$n$ Le Cam layer of Prop.~\ref{prop:knowcap-finiten} carries over
with $\Phi\to F$).
\end{theorem}
\begin{proof}
The converse and achievability are exactly Theorem~\ref{thm:knowcap} with $\Phi\to F$:
the admissible set is the interval $[p_S-\varepsilon,p_S+\varepsilon]$ in the mass
coordinate $u=F(\theta_T-\mu)$, the flip is at the band-median mass
$u_0=\tfrac12[F(-\mu)+F(a-\mu)]$, and $K\gtrless1\iff u_0$ lies outside/inside the
admissible band; identical-$Q_X$ gives $\TV=0$. The only place a hypothesis is used is
the claim that $\theta_0(\mu)$ is monotone in $\mu$ (needed for $K(\mu)$ to be a single
threshold rather than a fragmented set): for an MLR/log-concave location family the map
$\mu\mapsto\theta_0(\mu)$ is strictly increasing because $F(a-\mu)$ and $F(-\mu)$ are
decreasing in $\mu$ and $F^{-1}$ is increasing, with no antimode to reverse the band
median's motion. \emph{(Validated: Laplace and logistic location families, $0$
mismatches over $6{,}000$ draws each, benefit identity matched to
$\le4.6\times10^{-4}$ MC error; \texttt{val\_knowability\_capacity\_general.py}
Stage~2a.)}
\end{proof}

\begin{theorem}[Stage 2, non-log-concave breaker; PROVEN]\label{thm:gen-breaker}
There is a location family for which the scalar capacity \emph{fails}. Take $q_0$ a
balanced mixture of two well-separated Gaussians (bimodal, hence not log-concave and
not MLR) and the band $D=(0,a)$ wide enough to span the antimode. Then although
$\Delta(\theta)$ in \eqref{eq:gen-fam-delta} is still monotone in $\theta$ (so the
concept-location flip is unique), the flip locus $\theta_0(\mu)$ is \textbf{non-monotone
in the unknown shift} $\mu$: as $\mu$ moves the band median across the antimode,
$\theta_0(\mu)$ collapses. Consequently the label-free identifiable set
$\{\mu:K(\mu)>1\}$ \textbf{fragments} into more components than in the unimodal case, and
no single scalar threshold $\tau$ separates the identifiable from the
non-identifiable regime.
\end{theorem}
\begin{proof}[Proof (by exhibited witness)]
The benefit identity \eqref{eq:gen-fam-delta} holds for the mixture CDF $F$, monotone
in $\theta$. The density $q_0$ has two modes (verified: two local maxima), so the
hypothesis of Theorem~\ref{thm:gen-mlr} fails. Tracing $\theta_0(\mu)$ over
$\mu\in[-3,3]$ exhibits two strict decreases (non-monotone events) where the band
median crosses the antimode, and the set $\{\mu:K(\mu)>1\}$ has $2$ components for the
mixture versus $1$ for the matched unimodal Gaussian baseline at the same
$(p_S,\varepsilon,a)$. Hence the achievability certificate ``estimate $\hat\mu$, read
off $\sign(\theta_S-\theta_0(\hat\mu))$'' is not governed by a single threshold in $K$.
\emph{(Validated: density modes $=2$, $\theta_0(\mu)$ non-monotone events $=2$,
components mixture$=2$ vs Gaussian$=1$; \texttt{val\_knowability\_capacity\_general.py}
Stage~2b.)}
\end{proof}

% ----------------------------------------------------------------------------
\subsection{Stage 3: a distribution-free converse, and the general obstruction}
\label{sec:gen-df}
% ----------------------------------------------------------------------------
The converse needs none of the above structure. It is a pure two-point Le Cam bound
that uses only the defining feature of the problem: the concept lives in the
unobserved $P(Y\mid X)$ channel, so any two concepts induce the \emph{same} observable
covariate law.

\begin{theorem}[Distribution-free sign-of-benefit converse; PROVEN]\label{thm:gen-df-conv}
Let $X$ take values in any measurable space, let $g$ and $f$ be any two fixed
classifiers, and let $\mathcal H$ be a class of admissible concepts $\eta:\,
\mathcal X\to[0,1]$. Suppose there exist $\eta^+,\eta^-\in\mathcal H$ that induce the
\textbf{same observable marginal} $Q_X$ (automatic under covariate shift, where the
concept does not affect $Q_X$) and have \textbf{opposite benefit sign},
$\sign\big(R_Q^{\eta^+}(g)-R_Q^{\eta^+}(f)\big)=-\sign\big(R_Q^{\eta^-}(g)-R_Q^{\eta^-}(f)\big)\neq0$.
Then for \emph{every} sample size $n$ and \emph{every} label-free decision rule
$\hat T$ that observes only $\{X_i\}_{i=1}^n\sim Q_X^{\otimes n}$,
\[
  \inf_{\hat T}\ \max_{\eta\in\{\eta^+,\eta^-\}}\ \Pr\big[\hat T\ \text{outputs the wrong }\sign\Delta\big]
  \;\ge\;\tfrac12\big(1-\TV(Q_X^{\otimes n},Q_X^{\otimes n})\big)\;=\;\tfrac12 .
\]
No distributional, dimensional, or parametric assumption is used.
\end{theorem}
\begin{proof}
The two worlds $\eta^\pm$ generate identical laws for the unlabeled sample
($Q_X^{\otimes n}$ in both), so the total variation between the observed data
distributions is $0$. A label-free sign decision is a test between the two worlds,
which have opposite correct answers; Le Cam's two-point inequality gives minimax error
$\ge\tfrac12(1-\TV)=\tfrac12$, attained by a constant guess.
\emph{(Validated across Gaussian-1d, Laplace-1d, and aligned Gaussian-3d: opposite
$\Delta$ signs with unlabeled $\TV=0$ and minimax lower bound $\tfrac12$;
\texttt{val\_knowability\_capacity\_general.py} Stage~3a.)}
\end{proof}

\noindent\textbf{The general obstruction (why achievability needs a condition).}
Theorem~\ref{thm:gen-df-conv} is fully general, but a matching \emph{achievability}
(a single computable $K$ with the property $K>1\Rightarrow$ identifiable) provably
cannot be. The achievability side requires the observer to (a) locate the
sign-flip set of $\Delta$ over the admissible class and (b) certify that the whole
admissible class lies on one side of it from $Q_X$ alone. Stage~1b shows step~(a)
already fails to be captured by an alignment-blind margin once the concept direction is
free (the flip locus depends on the Mahalanobis geometry between concept and rule), and
Stage~2b shows that, even in one dimension, when the family is not MLR/unimodal the
flip locus is non-monotone in the unknown shift, so step~(b) has no single-threshold
certificate. We record the general statement as a conjecture with the precise missing
piece.

\begin{conjecture}[General knowability--capacity; OPEN, with the exact gap]
\label{conj:gen-capacity}
Fix observables $(Q_X\text{-family},g,f)$ and an admissible concept class $\mathcal H$
defined by a mass-drift budget $\varepsilon$. Suppose a \textbf{regularity/identifiability
condition} holds: \textnormal{(R1)} the benefit $\Delta$ has, over $\mathcal H$, a
\emph{unique} flip locus expressible in the observable mass coordinate; and
\textnormal{(R2)} that flip locus is a \emph{monotone, observable} function of the
unknown nuisance (e.g.\ the shift) --- guaranteed by MLR/log-concavity in the location
case and by Mahalanobis-alignment in the multivariate case. Then the scalar
\[
  K\;=\;\frac{\operatorname{dist}_{\text{mass}}\!\big(\text{calibrated boundary},\ \text{worst-case flip locus over }\mathcal H\big)}{\varepsilon}
\]
is a capacity at $\tau=1$: $K>1\iff$ the sign of $\Delta$ is label-free identifiable,
and the finite-$n$ price is an $O(1/\sqrt n)$ Le Cam layer around the threshold
\emph{location}. The \textbf{open part} is precisely the removal (or sharp
characterization) of \textnormal{(R1)--(R2)}: without them the flip set need not be a
single observable threshold, the worst-case-flip distance need not be computable from
$Q_X$, and no scalar $\tau$ is simultaneously necessary and sufficient. Stages~1--2
prove the conjecture in its two leading regularity classes and exhibit explicit
breakers (free-direction tilt; non-log-concave bimodality) on the boundary of each.
\end{conjecture}

\twocolumn
\begin{remark}[Honest scope of the generalization]\label{rmk:gen-scope}
This section reaches a \emph{strong but bounded} generalization, \textbf{not} a
paradigm-level general capacity. What is \emph{proven}: (i) an \emph{exact}
$d$-dimensional lift of $\tau=1$ in Mahalanobis geometry for aligned concepts
(Thm~\ref{thm:gen-aligned}); (ii) the \emph{exact} obstruction and worst-case capacity
$K_{\mathrm{free}}$ for free halfspace concepts (Thm~\ref{thm:gen-free}); (iii) a
matched $\tau=1$ for MLR/log-concave location families (Thm~\ref{thm:gen-mlr}) with an
explicit non-log-concave breaker (Thm~\ref{thm:gen-breaker}); and (iv) a fully
\emph{distribution-free converse} in any dimension and any family
(Thm~\ref{thm:gen-df-conv}). What is \emph{not} proven, and is stated as
Conjecture~\ref{conj:gen-capacity}: a single computable capacity \emph{without} a
regularity/identifiability condition. The wall is sharp and named: a scalar $\tau$ is a
capacity \emph{iff} the benefit's flip locus is a unique, monotone, observable
threshold; alignment (multivariate) and MLR/unimodality (location) are the two
sufficient regularity conditions we can currently prove, and each has a concrete
counterexample just outside it. The value of the result is therefore (a) the converse
is essentially universal, and (b) the achievability is pinned to an explicit,
checkable regularity hypothesis whose failure modes are exhibited --- a precise map of
how far the $1$-D capacity travels and exactly where it stops.
\end{remark}
